%Specify document class
\documentclass[12pt,twoside,openany]{book}

%Packages to include
\usepackage[utf8]{inputenc}
\usepackage{setspace}
\usepackage{times}
\usepackage[a4paper,width=150mm,top=30mm,bottom=1in,bindingoffset=6mm]{geometry}
\usepackage{graphicx,wrapfig,lipsum}
\usepackage{bm}
\usepackage{indentfirst}
\usepackage{verbatim}
\usepackage{float}
\usepackage[shortlabels]{enumitem}
\usepackage{subfigure}
\usepackage{titlesec}
\usepackage[font={footnotesize}]{caption}
\usepackage{amsmath}
\usepackage{amssymb}
\usepackage[subfigure,titles]{tocloft}
\usepackage{cite}
\usepackage{url}
\usepackage{xfrac}
\usepackage[nottoc,numbib]{tocbibind}
\usepackage[table,xcdraw]{xcolor}
\usepackage{siunitx}
\usepackage{lineno}
\settocbibname{References}
\usepackage[%  
    colorlinks=true,
    pdfborder={0 0 0},
    linkcolor=blue,
    citecolor=blue,
    urlcolor=blue
]{hyperref}
\usepackage[nameinlink]{cleveref}

\linenumbers

%Set images folder
\graphicspath{ {images/} }

%Set Header and Footer for all pages
\usepackage{fancyhdr}

\pagestyle{fancy}
\fancyhead{}
\fancyhead[LO]{\nouppercase{\textbf{\leftmark}}}
\fancyhead[RO,LE]{\textbf{\thepage}}
\fancyhead[RE]{\nouppercase{\textbf{\rightmark}}}
\fancyfoot{}
\renewcommand{\chaptermark}[1]{%
\markboth{#1}{}}

%Set Title Page Background
\usepackage[pages=some]{background}
\backgroundsetup{
scale=1,
color=black,
opacity=0.1,
angle=0,
contents={\includegraphics[width=\textwidth]{university.png}}
}

\setlist[itemize]{topsep=\parskip}

%\setlength\abovecaptionskip{-5ex}
%\setlength{\textfloatsep}{0pt plus 1.0pt minus 2.0pt}
\setlength\belowcaptionskip{-3ex}

\usepackage[labelfont=bf]{caption}
\captionsetup{labelfont=bf}

%Eliminate extra page after title and table of contents
\let\cleardoublepage=\clearpage

%Make signature and date lines for title page
\newcommand*{\SignatureAndDate}[1]{
    \par\noindent\makebox[3.0in]{\hrulefill} \hfill\makebox[2.0in]{\hrulefill}
    \par\noindent\makebox[2.5in][l]{#1}      \hfill\makebox[2.0in][l]{Date}
}

%for list of figures and tables
\newlength{\mylen}

\renewcommand{\cftfigpresnum}{\figurename\enspace}
\renewcommand{\cftfigaftersnum}{:}
\settowidth{\mylen}{\cftfigpresnum\cftfigaftersnum}
\addtolength{\cftfignumwidth}{\mylen}

%Main Body
\begin{document}

%Title Page
\begin{titlepage}
    \begin{center}
    
      % \ThisCenterWallPaper{1}{university.png}
       \BgThispage
        \vspace*{1cm}
      \Large  
        \textbf{Measurement of the Z Invisible background for Stop Quark Searches in all Hadronic Channel at \bm{$\sqrt{s} = 13$} TeV}
        
        \normalsize
        \vspace{0.5cm}
        by
        
        \vspace{0.25cm}
        
        Andr\'es J. Abreu Nazario
        
        \vspace{0.5cm}	
        
        A thesis presented for the degree of\\
        \vspace{0.5cm}
        MASTER OF SCIENCE\\
        \vspace{0.5cm}
        in\\
        \vspace{0.5cm}
        Physics\\
        \vspace{0.5cm}
        UNIVERSITY OF PUERTO RICO\\
        MAYAG\"UEZ CAMPUS\\
        2018
        \vspace{0.8cm}
        
        \end{center}
        
        Approved by:
        
        \vspace{1.0cm}

\SignatureAndDate{Sudhir Malik, Ph.D.}
\\President, Graduate Committee

\vspace{1.0cm}

\SignatureAndDate{H\'ector M\'endez, Ph.D.}
\\Member, Graduate Committee

\vspace{1.0cm}

\SignatureAndDate{Samuel Santana Col\'on, Ph.D.}
\\Member, Graduate Committee

\vspace{1.0cm}

\SignatureAndDate{Rafael A. R\'amos, Ph.D.}
\\Chairperson of the Department



\end{titlepage}

\frontmatter

\doublespace
\chapter{Abstract}
The Data Quality Monitoring (DQM) of CMS is a key asset to deliver high-quality data for physics analysis and it is used both in the online and offline environment. The current paradigm of the quality assessment is labor intensive and it is based on the scrutiny of a large number of histograms by detector experts comparing them with a reference. This project aims at applying recent progress in Machine Learning techniques to the automation of the DQM scrutiny. In particular the use of convolutional neural networks to spot problems in the acquired data is presented with particular attention to semi-supervised models (e.g. autoencoders) to define a classification strategy that doesn’t assume previous knowledge of failure modes. Real data from the hadron calorimeter of CMS are used to demonstrate the effectiveness of the proposed approach. 

\vspace*{1cm}

\textit{Keywords}:  [DQM, online, offline, Machine Learning ]

\chapter{Acknowledgments}

I wish to thank United States State Department and University of Michigan for providing the opportunity to work abroad at CERN during the 2018 Winter Semester. I also wish to thank CERN staff, CMS Experiment , Texas Tech University, and the University of Puerto Rico at Mayagüez, with special thanks to Dr. Federico de Guio for his local mentorship and Dr. Nural Akchurin, and Dr. Sudhir Malik for their guidance. A very special thanks to Dr. Jean Krisch for accepting me for this great research opportunity and Dr. Steven Goldfarb for being a wonderful host and overall local guidance at CERN.
\singlespace

\listoffigures

\tableofcontents


\doublespace
\mainmatter
\chapter{Introduction}
\input{chapters/introduction}

\chapter{The CMS Experiment}\label{LHCCMS}
\input{chapters/chapter02}

\chapter{Data Collection and Data Quality Monitoring }\label{DQMchapter}
\input{chapters/chapter03}

\chapter{What is Machine Learning?}\label{ML}
\input{chapters/chapter04}

%\chapter{Search for SUSY in the 0-Lepton Final State}\label{AnalysisChap}
%\input{chapters/chapter04}
%
%\chapter{Estimation of the Z$\rightarrow \nu\bar{\nu}$ Background}\label{estimation}
%\input{chapters/chapter05}
%
%\chapter{Conclusion}\label{conclusion}
%\input{chapters/conclusion}

\appendix
\chapter{Appendix Title}
%\input{chapters/appendix}

\bibliography{references}{}
\bibliographystyle{ieeetr}

\end{document}
